% GENERATED FILE. Edit canonical lessons, then run npm run generate:books.
% canonical-source-hash: fnv1a64:b05d341e30dfd8b7
% canonical-lessons: HI-C492-malis, HI-C492-svasth, HI-C492-canga, HI-C492-ayurved, HI-C492-sphurti

\chapter{Getting Better}
\label{ch:hi-getting-better}
\hlchaptermodality{492}

\begin{chapteropening}
\textbf{By the end of this chapter:} \emph{I can say that someone is getting better.}

Complete the last lesson of chapter 492: I can say that someone is getting better.
\end{chapteropening}

\section[\texorpdfstring{\hi{मालिश} (māliś)}{मालिश (māliś)}]{\hi{मालिश} (māliś) — a massage; a rub}

\label{lesson:HI-C492-malis}



\begin{quote}
Before the new one: say the Hindi for to have someone sent for, then the Hindi for a life.
\end{quote}

\subsection*{You'll want to know: \hi{मालिश}}
\textbf{\hi{मालिश}} — \emph{māliś} — \textquotedblleft{}a massage; a rub\textquotedblright{}.

\begin{quote}
\hi{रोज़} \hi{की} \hi{मालिश} \hi{से} \hi{दादी} \hi{के} \hi{घुटनों} \hi{का} \hi{दर्द} \hi{कम} \hi{हो} \hi{गया।}
\end{quote}

\begin{quote}
\emph{roz kī māliś se dādī ke ghuṭnoṁ kā dard kam ho gayā.}
\end{quote}

\begin{quote}
With a daily massage, the pain in my grandmother's knees eased.
\end{quote}

Feminine.

\begin{cousinweb}
\textbf{\hi{मालिश}} (\emph{māliś}) is usually traced to Persian \emph{mālish}, \textquotedblleft{}rubbing\textquotedblright{}, the noun of the Persian verb for \textquotedblleft{}to rub\textquotedblright{}. \emph{māliś karnā} is to massage; \emph{malnā}, \textquotedblleft{}to rub in\textquotedblright{}, is a separate, inherited word.
\end{cousinweb}

\begin{grammarlens}[title={What you've built}]
One more word for describing what you lived and felt.
\end{grammarlens}

\subsection*{Your turn}
\begin{itemize}
\raggedright
  \item \emph{Say it:} \emph{māliś}
  \item \emph{Say it:} \emph{māliś}, once more
  \item \emph{Say it:} \emph{roz kī māliś se dādī ke ghuṭnoṁ kā dard kam ho gayā.} — With a daily massage, the pain in my grandmother's knees eased.
  \item \emph{Recall:} say the Hindi for to have someone sent for, then the Hindi for a life, then say \emph{māliś} again
\end{itemize}

\begin{tcolorbox}[breakable,colback=teal!4,colframe=teal!35!black,title={Before you move on}]
What does \hi{मालिश} mean? (\textquotedblleft{}A massage; a rub\textquotedblright{}.) Say it once more.
\end{tcolorbox}
\section[\texorpdfstring{\hi{स्वस्थ} (svasth)}{स्वस्थ (svasth)}]{\hi{स्वस्थ} (svasth) — healthy; well}

\label{lesson:HI-C492-svasth}



\begin{quote}
Before the new one: say the Hindi for a life, then the Hindi for a massage.
\end{quote}

\subsection*{You'll want to know: \hi{स्वस्थ}}
\textbf{\hi{स्वस्थ}} — \emph{svasth} — \textquotedblleft{}healthy; well\textquotedblright{}.

\begin{quote}
\hi{रोज़} \hi{दौड़कर} \hi{अरुण} \hi{स्वस्थ} \hi{रहता} \hi{है।}
\end{quote}

\begin{quote}
\emph{roz dauṛkar Aruṇ svasth rahtā hai.}
\end{quote}

\begin{quote}
Arun keeps healthy by running every day.
\end{quote}

\begin{cousinweb}
\textbf{\hi{स्वस्थ}} (\emph{svasth}) is a tatsama, borrowed from Sanskrit \emph{svastha}, \textquotedblleft{}self-abiding, being in one's natural state\textquotedblright{}, and so \textquotedblleft{}well, healthy\textquotedblright{} (Monier-Williams), built on \emph{sva}, \textquotedblleft{}self\textquotedblright{}, and \emph{stha}, \textquotedblleft{}standing\textquotedblright{}.
\end{cousinweb}

\begin{grammarlens}[title={What you've built}]
One more word for describing what you lived and felt.
\end{grammarlens}

\subsection*{Your turn}
\begin{itemize}
\raggedright
  \item \emph{Say it:} \emph{svasth}
  \item \emph{Say it:} \emph{svasth}, once more
  \item \emph{Say it:} \emph{roz dauṛkar Aruṇ svasth rahtā hai.} — Arun keeps healthy by running every day.
  \item \emph{Recall:} say the Hindi for a life, then the Hindi for a massage, then say \emph{svasth} again
\end{itemize}

\begin{tcolorbox}[breakable,colback=teal!4,colframe=teal!35!black,title={Before you move on}]
What does \hi{स्वस्थ} mean? (\textquotedblleft{}Healthy; well\textquotedblright{}.) Say it once more.
\end{tcolorbox}
\section[\texorpdfstring{\hi{चंगा} (caṅgā)}{चंगा (caṅgā)}]{\hi{चंगा} (caṅgā) — well again; cured}

\label{lesson:HI-C492-canga}



\begin{quote}
Before the new one: say the Hindi for a massage, then the Hindi for healthy.
\end{quote}

\subsection*{You'll want to know: \hi{चंगा}}
\textbf{\hi{चंगा}} — \emph{caṅgā} — \textquotedblleft{}well again; cured\textquotedblright{}.

\begin{quote}
\hi{आराम} \hi{से} \hi{तुम} \hi{जल्दी} \hi{चंगे} \hi{हो} \hi{जाओगे।}
\end{quote}

\begin{quote}
\emph{ārām se tum jaldī caṅge ho jāoge.}
\end{quote}

\begin{quote}
With rest you'll soon be well again.
\end{quote}

\begin{cousinweb}
\textbf{\hi{चंगा}} (\emph{caṅgā}) is a tadbhava: CDIAL lists Hindi \emph{caṅgā}, \textquotedblleft{}good, sound, well\textquotedblright{}, under Sanskrit \emph{caṅga} (CDIAL 4564).
\end{cousinweb}

\begin{grammarlens}[title={What you've built}]
One more word for describing what you lived and felt.
\end{grammarlens}

\subsection*{Your turn}
\begin{itemize}
\raggedright
  \item \emph{Say it:} \emph{caṅgā}
  \item \emph{Say it:} \emph{caṅgā}, once more
  \item \emph{Say it:} \emph{ārām se tum jaldī caṅge ho jāoge.} — With rest you'll soon be well again.
  \item \emph{Recall:} say the Hindi for a massage, then the Hindi for healthy, then say \emph{caṅgā} again
\end{itemize}

\begin{tcolorbox}[breakable,colback=teal!4,colframe=teal!35!black,title={Before you move on}]
What does \hi{चंगा} mean? (\textquotedblleft{}Well again; cured\textquotedblright{}.) Say it once more.
\end{tcolorbox}
\section[\texorpdfstring{\hi{आयुर्वेद} (āyurved)}{आयुर्वेद (āyurved)}]{\hi{आयुर्वेद} (āyurved) — ayurveda; traditional Indian medicine}

\label{lesson:HI-C492-ayurved}



\begin{quote}
Before the new one: say the Hindi for healthy, then the Hindi for well again.
\end{quote}

\subsection*{You'll want to know: \hi{आयुर्वेद}}
\textbf{\hi{आयुर्वेद}} — \emph{āyurved} — \textquotedblleft{}ayurveda; traditional Indian medicine\textquotedblright{}.

\begin{quote}
\hi{आयुर्वेद} \hi{की} \hi{दवा} \hi{से} \hi{दादा} \hi{जल्दी} \hi{ठीक} \hi{हो} \hi{गए।}
\end{quote}

\begin{quote}
\emph{āyurved kī davā se dādā jaldī ṭhīk ho gae.}
\end{quote}

\begin{quote}
With ayurvedic medicine, my grandfather soon got well.
\end{quote}

Masculine.

\begin{cousinweb}
\textbf{\hi{आयुर्वेद}} (\emph{āyurved}) is a tatsama, borrowed from Sanskrit \emph{āyurveda}, which Monier-Williams explains as the learning of health and healing (Monier-Williams), built on \emph{āyus}, \textquotedblleft{}life\textquotedblright{}, and \emph{veda}, \textquotedblleft{}knowledge\textquotedblright{}. It names India's traditional medicine and its herbal remedies.
\end{cousinweb}

\begin{grammarlens}[title={What you've built}]
One more word for describing what you lived and felt.
\end{grammarlens}

\subsection*{Your turn}
\begin{itemize}
\raggedright
  \item \emph{Say it:} \emph{āyurved}
  \item \emph{Say it:} \emph{āyurved}, once more
  \item \emph{Say it:} \emph{āyurved kī davā se dādā jaldī ṭhīk ho gae.} — With ayurvedic medicine, my grandfather soon got well.
  \item \emph{Recall:} say the Hindi for healthy, then the Hindi for well again, then say \emph{āyurved} again
\end{itemize}

\begin{tcolorbox}[breakable,colback=teal!4,colframe=teal!35!black,title={Before you move on}]
What does \hi{आयुर्वेद} mean? (\textquotedblleft{}Ayurveda; traditional Indian medicine\textquotedblright{}.) Say it once more.
\end{tcolorbox}
\section[\texorpdfstring{\hi{स्फूर्ति} (sphūrti)}{स्फूर्ति (sphūrti)}]{\hi{स्फूर्ति} (sphūrti) — vigour; fresh energy}

\label{lesson:HI-C492-sphurti}



\begin{quote}
Before the new one: say the Hindi for well again, then the Hindi for ayurveda.
\end{quote}

\subsection*{You'll want to know: \hi{स्फूर्ति}}
\textbf{\hi{स्फूर्ति}} — \emph{sphūrti} — \textquotedblleft{}vigour; fresh energy\textquotedblright{}.

\begin{quote}
\hi{सुबह} \hi{की} \hi{सैर} \hi{के} \hi{बाद} \hi{दादी} \hi{में} \hi{नई} \hi{स्फूर्ति} \hi{थी।}
\end{quote}

\begin{quote}
\emph{subah kī sair ke bād dādī meṁ naī sphūrti thī.}
\end{quote}

\begin{quote}
After her morning walk my grandmother felt full of fresh energy.
\end{quote}

Feminine.

\begin{cousinweb}
\textbf{\hi{स्फूर्ति}} (\emph{sphūrti}) is a tatsama, borrowed from Sanskrit \emph{sphūrti}, \textquotedblleft{}quivering, throbbing\textquotedblright{} (Monier-Williams); Hindi uses it for the lively energy one feels after rest, a sense the dictionary does not give. It is the learned cousin of \emph{phurtī}.
\end{cousinweb}

\begin{grammarlens}[title={What you've built}]
One more word for describing what you lived and felt.
\end{grammarlens}

\subsection*{Your turn}
\begin{itemize}
\raggedright
  \item \emph{Say it:} \emph{sphūrti}
  \item \emph{Say it:} \emph{sphūrti}, once more
  \item \emph{Say it:} \emph{subah kī sair ke bād dādī meṁ naī sphūrti thī.} — After her morning walk my grandmother felt full of fresh energy.
  \item \emph{Recall:} say the Hindi for well again, then the Hindi for ayurveda, then say \emph{sphūrti} again
\end{itemize}

\begin{tcolorbox}[breakable,colback=teal!4,colframe=teal!35!black,title={Before you move on}]
What does \hi{स्फूर्ति} mean? (\textquotedblleft{}Vigour; fresh energy\textquotedblright{}.) Say it once more.
\end{tcolorbox}
